\documentclass[12pt]{article}
\usepackage[utf8]{inputenc}
\usepackage[T1]{fontenc}
\usepackage{hyperref}
\title{Tobo Language Resources for K-12 Curriculum Development}
\author{Compiled by Claude Sonnet 5 (claude-sonnet-5) for Andr\'as Kornai}
\date{}
\begin{document}
\maketitle
\section{Introduction}
Tobo (ISO 639-3 tbv; Glottocode tobo1251 per transnewguinea.org) is a Trans-New Guinea (Finisterre-Huon) language of Morobe Province, Papua New Guinea. It is closely related to Kube (kgf) and the two are reported to be mutually intelligible. Speaker estimates range from about 2,230 (1980 census) to about 3,000. Little material exists online.
\section{Primary Educational Resources}
No primers, textbooks or school programmes were found.
\section{Linguistic and Reference Materials}
\begin{itemize}
\item Wikipedia, Tobo-Kube language. \url{https://en.wikipedia.org/wiki/Tobo-Kube_language}
\item Tobo wordlist on the Trans New Guinea database (based on McElhanon 1967 and McElhanon and Voorhoeve 1970). \url{https://transnewguinea.org/language/tobo}
\item Numeral data. \url{https://lingweb.eva.mpg.de/channumerals/Tobo.htm}
\end{itemize}
\section{Religious and Community Resources}
Joshua Project people-group page for the Tobo (notes limited published detail). \url{https://joshuaproject.net/people_groups/15524}
\section{Digital and Audio Resources}
No audio or video found in the searches.
\section{Curriculum Development Context}
No grammar, dictionary, or literacy materials were found, so a curriculum would start from scratch, with community leadership and fieldwork. The close relation to Kube may allow partly shared materials, but this needs verification with speakers.
\end{document}